\chapter{Multiplier des nombres relatifs}\label{ch:g8:negprod}

L'année 7 a additionné et soustrait des \omterm{def:g7:negatives:def}{nombres relatifs}
(\cref{ch:g7:negatives})~; il reste à les multiplier et les diviser.
Une question domine le chapitre~: pourquoi «~moins fois moins~» doit-il
être plus~? On donne une raison, pas seulement une règle.

\section{Les règles de signes}

\begin{theorem}[Signe d'un produit]\label{thm:g8:negprod:rules}
Le \omterm{def:g2:mult:def}{produit} de deux \omterm{def:g7:negatives:def}{nombres relatifs} a pour distance à \omterm{def:g1:counting:zero}{zéro} le \omterm{def:g2:mult:def}{produit}
des distances, et son signe est donné par~:
\begin{center}
\renewcommand{\arraystretch}{1.3}
\begin{tabular}{c|cc}
$\times$ & $+$ & $-$ \\
\hline
$+$ & $+$ & $-$ \\
$-$ & $-$ & $+$ \\
\end{tabular}
\end{center}
\emph{Mêmes signes~: \omterm{def:g2:mult:def}{produit} positif. Signes contraires~: \omterm{def:g2:mult:def}{produit}
négatif.} Le même tableau régit les \omterm{def:g3:division:remainder}{quotients}.
\end{theorem}

\begin{proof}[Pourquoi $(-1) \times (-1) = +1$]
D'abord, $3 \times (-4)$ signifie $(-4) + (-4) + (-4) = -12$~: un
positif fois un négatif est négatif. Observer maintenant le motif~:
\[
3 \times (-4) = -12, \quad
2 \times (-4) = -8, \quad
1 \times (-4) = -4, \quad
0 \times (-4) = 0 .
\]
À chaque étape, le premier facteur baisse de $1$ et le résultat
\emph{monte} de $4$. En poursuivant le motif d'un pas~:
\[
(-1) \times (-4) = 0 + 4 = +4 .
\]
Tout autre choix casserait la régularité de l'arithmétique
(distributivité). Donc moins fois moins doit être plus.
\end{proof}

\begin{omfigure}
\begin{tikzpicture}
\begin{axis}[omaxis,
  width=7.8cm, height=5.4cm,
  xmin=-2.6, xmax=3.6, ymin=-14, ymax=10,
  xlabel={$k$}, ylabel={$k \times (-4)$},
  xtick={-2,-1,1,2,3}, ytick={-12,-8,-4,4,8}]
  \addplot[omDef, thick, domain=-2.3:3.3] {-4*x};
  \addplot[omDef, only marks, mark=*, mark size=1.6pt] coordinates
    {(3,-12) (2,-8) (1,-4) (0,0) (-1,4) (-2,8)};
\end{axis}
\end{tikzpicture}

{\small L'argument du motif en image~: les points
$(k,\ k \times (-4))$ s'alignent, et prolonger la droite au-delà de
$k = 0$ force $(-1)\times(-4) = +4$ et $(-2)\times(-4) = +8$.}
\end{omfigure}

\begin{example}\label{ex:g8:negprod:products}
\begin{align*}
(-7) \times (+6) &= -42 && \text{(signes contraires)} \\
(-5) \times (-8) &= +40 && \text{(mêmes signes)} \\
(+2.5) \times (-4) &= -10 \\
(-36) \div (-9) &= +4 \\
(+15) \div (-2) &= -7.5 .
\end{align*}
\end{example}

\section{Produits de plusieurs facteurs}

\begin{proposition}[Signe d'un long produit]\label{prop:g8:negprod:many}
Le signe d'un \omterm{def:g2:mult:def}{produit} de plusieurs facteurs non nuls ne dépend que du
nombre de facteurs négatifs~:
\begin{itemize}
  \item un nombre \emph{\omterm{def:g3:numbers:evenodd}{pair}} de facteurs négatifs donne un \omterm{def:g2:mult:def}{produit}
        positif~;
  \item un nombre \emph{\omterm{def:g3:numbers:evenodd}{impair}} donne un \omterm{def:g2:mult:def}{produit} négatif.
\end{itemize}
\end{proposition}

\begin{proof}
Les facteurs négatifs se regroupent par paires, et chaque paire
contribue un signe positif (\cref{thm:g8:negprod:rules})~; un décompte
\omterm{def:g3:numbers:evenodd}{impair} laisse un facteur négatif non apparié, qui rend le \omterm{def:g2:mult:def}{produit}
entier négatif.
\end{proof}

\begin{example}\label{ex:g8:negprod:many}
$(-2) \times (+3) \times (-5) \times (-1)$~: trois facteurs négatifs
(\omterm{def:g3:numbers:evenodd}{impair}), donc le \omterm{def:g2:mult:def}{produit} est négatif~; distances~:
$2 \times 3 \times 5 \times 1 = 30$. Résultat~: $-30$. Décider le
signe \emph{d'abord}, puis multiplier les distances --- deux tâches
faciles au lieu d'une sujette aux erreurs.
\end{example}

\begin{example}[Puissances de négatifs]\label{ex:g8:negprod:powers}
$(-2)^2 = (-2) \times (-2) = +4$, et
$(-2)^3 = (-2)^2 \times (-2) = -8$~: les exposants \omterm{def:g3:numbers:evenodd}{pairs} donnent des
valeurs positives, les exposants \omterm{def:g3:numbers:evenodd}{impairs} gardent le signe. Attention
à la notation~: $(-3)^2 = 9$ mais $-3^2 = -(3 \times 3) = -9$.
\end{example}

\section{Calculer avec les quatre opérations}

\begin{method}[Calcul sûr avec des nombres relatifs]\label{met:g8:negprod:compute}
\begin{enumerate}
  \item Respecter les priorités du \cref{ch:g7:priorities}~:
        parenthèses, puis $\times$ et $\div$, puis $+$ et $-$~;
  \item à chaque multiplication ou division, trouver le \emph{signe
        d'abord}, puis la distance~;
  \item réécrire une étape par ligne.
\end{enumerate}
\end{method}

\begin{example}\label{ex:g8:negprod:compute}
\begin{align*}
5 - 3 \times (-4)
&= 5 - (-12) && \text{(produit d'abord~: signes contraires)}\\
&= 5 + 12 = 17 .
\end{align*}
\begin{align*}
\frac{(-6) \times (-10)}{-4}
&= \frac{+60}{-4} && \text{(numérateur~: mêmes signes)}\\
&= -15 && \text{(quotient~: signes contraires).}
\end{align*}
\end{example}

\begin{example}[Substituer des valeurs négatives]\label{ex:g8:negprod:substitute}
Évaluer $E = 3x^2 - 2x$ pour $x = -5$, parenthèses autour de la
valeur~:
\[
E = 3 \times (-5)^2 - 2 \times (-5)
= 3 \times 25 - (-10)
= 75 + 10 = 85 .
\]
Les parenthèses de $(-5)^2$ sont essentielles~: sans elles le carré
ne s'appliquerait qu'à $5$.
\end{example}

\section{Exercices}

\begin{exercise}[$\star$]\label{exo:g8:negprod:1}
Calculer~:
\[
(-8) \times (+7), \qquad
(-9) \times (-6), \qquad
(+12) \times (-0.5), \qquad
(-1) \times (-1) \times (-1).
\]
\end{exercise}

\begin{exercise}[$\star$]\label{exo:g8:negprod:2}
Calculer~:
\[
(-63) \div (+9), \qquad
(-8) \div (-16), \qquad
\frac{45}{-5}, \qquad
\frac{-7.2}{-8} .
\]
\end{exercise}

\begin{exercise}[$\star$]\label{exo:g8:negprod:3}
Donner seulement le \emph{signe} de chaque \omterm{def:g2:mult:def}{produit}, sans le
calculer~:
\[
(-3) \times (-7) \times (+2) \times (-5); \qquad
(-1)^{10}; \qquad
(-2)^{7}; \qquad
(-4) \times 0 \times (-6).
\]
\end{exercise}

\begin{exercise}[$\star$]\label{exo:g8:negprod:4}
Calculer~:
\[
(-2)^4, \qquad
-2^4, \qquad
(-3)^3, \qquad
(-1)^{2026} .
\]
\end{exercise}

\begin{exercise}[$\star$]\label{exo:g8:negprod:5}
Calculer étape par étape, en respectant les priorités~:
\[
7 + 2 \times (-6), \qquad
(-4) \times (5 - 9), \qquad
18 \div (-3) - (-2) .
\]
\end{exercise}

\begin{exercise}[$\star$]\label{exo:g8:negprod:6}
Calculer~:
\[
\frac{(-5) \times (+8)}{-4}, \qquad
\frac{(-3) \times (-6)}{(-2) \times (+9)} .
\]
\end{exercise}

\begin{exercise}[$\star$]\label{exo:g8:negprod:7}
Évaluer pour $x = -3$~:
\[
5x, \qquad x^2, \qquad -x^2, \qquad 2x^2 + 4x, \qquad (2x)^2 .
\]
\end{exercise}

\begin{exercise}[$\star\star$]\label{exo:g8:negprod:8}
Compléter chaque égalité~:
\[
(-6) \times \;?\; = 42, \qquad
\;?\; \div (-4) = -2.5, \qquad
(-5) \times \;?\; = -1 .
\]
\end{exercise}

\begin{exercise}[$\star\star$]\label{exo:g8:negprod:9}
Vrai ou faux~? Justifier ou donner un contre-exemple.
\begin{enumerate}
  \item Le \omterm{def:g2:mult:def}{produit} de deux \omterm{def:g7:negatives:def}{nombres relatifs} est toujours au moins
        aussi grand que leur \omterm{def:g1:addition:def}{somme}.
  \item Le carré d'un \omterm{def:g7:negatives:def}{nombre relatif} n'est jamais négatif.
  \item Si un \omterm{def:g2:mult:def}{produit} de trois facteurs est positif, les trois
        facteurs sont positifs.
\end{enumerate}
\end{exercise}

\begin{exercise}[$\star\star$]\label{exo:g8:negprod:10}
Chaque mauvaise réponse à un quiz compte $-3$ points, chaque bonne
réponse $+5$. Zoé a répondu à $20$ questions et a marqué $36$ points.
Combien de réponses étaient justes~? (Essayer des valeurs, ou mettre
en place le calcul $5r - 3(20 - r) = 36$.)
\end{exercise}

\begin{exercise}[$\star\star\star$]\label{exo:g8:negprod:11}
En utilisant la distributivité (comme dans la preuve de
\cref{thm:g8:negprod:rules}), développer et justifier chaque étape de
\[
0 = (-1) \times 0 = (-1) \times \bigl(1 + (-1)\bigr)
= (-1) \times 1 + (-1) \times (-1),
\]
et conclure que $(-1) \times (-1)$ doit valoir $+1$.
\end{exercise}

\section{Problème~: pourquoi moins par moins donne plus}

\begin{problem}[{Devoir maison --- les règles des signes déduites de
la distributivité, et la somme alternée $1 - 2 + 3 - 4 + \dots$}]
\label{pb:g8:negprod:1}
La preuve du~\cref{thm:g8:negprod:rules} a prolongé une régularité et
affirmé que «~tout autre choix casserait la distributivité~». Ce
problème rend l'affirmation exacte~: en partant de la seule
distributivité (\cref{thm:g7:priorities:distributivity}), on va
\emph{démontrer} les règles des signes --- sans dessin, sans
régularité, comme le fait l'algèbre --- puis les mettre au travail
sur les puissances de $-1$ et sur une \omterm{def:g1:addition:def}{somme} célèbre aux signes
alternés. Partout, $a$ et $b$ désignent des \omterm{def:g7:negatives:def}{nombres relatifs}~; on
admet les règles d'addition du \cref{ch:g7:negatives}, le fait que
$1 \times a = a$, et qu'un \omterm{def:g2:mult:def}{produit} se calcule dans n'importe quel
ordre (la distributivité s'applique donc des deux côtés d'un
\omterm{def:g2:mult:def}{produit}).

\medskip
\noindent\textbf{Partie I --- Les règles des signes sont des
théorèmes.}
\begin{enumerate}
  \item Démontrer que $0 \times a = 0$ pour tout \omterm{def:g7:negatives:def}{nombre relatif} $a$.
        (Écrire $0 = 0 + 0$, développer $(0 + 0) \times a$, et se
        demander quel nombre, ajouté à $0 \times a$, redonne
        $0 \times a$.)
  \item Développer $\bigl(1 + (-1)\bigr) \times a$ pour montrer que
        $(-1) \times a + a = 0$, et en déduire que
        \[
        (-1) \times a = -a :
        \]
        multiplier par $-1$ donne l'opposé. Comparer avec
        \cref{exo:g8:negprod:11}, qui est le cas $a = -1$.
  \item En déduire que $(-a) \times b = -(a \times b)$~: un signe
        moins sort d'un \omterm{def:g2:mult:def}{produit} sans changer.
  \item En déduire que $(-a) \times (-b) = a \times b$~: deux signes
        moins s'annulent.
  \item Tout \omterm{def:g7:negatives:def}{nombre relatif} est sa distance à \omterm{def:g1:counting:zero}{zéro} précédée d'un
        signe~: $a = +d$ ou $a = -d$. Expliquer comment les questions
        3 et 4 démontrent les quatre cas du tableau des signes du
        \cref{thm:g8:negprod:rules}, distances comprises.
\end{enumerate}

\medskip
\noindent\textbf{Partie II --- Les puissances de $-1$.}
\begin{enumerate}[resume]
  \item Calculer $(-1)^2$, $(-1)^3$, $(-1)^4$, $(-1)^5$, puis donner,
        en le justifiant par \cref{prop:g8:negprod:many}, la valeur
        de $(-1)^n$ pour tout entier $n \geq 1$.
  \item Donner (sans calculer aucune distance) le signe de
        \[
        (-1) \times (-2) \times (-3) \times \dots \times (-10),
        \]
        et écrire la distance à \omterm{def:g1:counting:zero}{zéro} de ce \omterm{def:g2:mult:def}{produit} sous la forme d'un
        \omterm{def:g2:mult:def}{produit} d'entiers, sans l'effectuer.
  \item Montrer que $(-a)^2 = a^2$ et que $(-a)^3 = -a^3$. Pour
        quels exposants le signe moins survit-il~?
  \item Utiliser les règles des signes pour montrer que le carré
        d'un \omterm{def:g7:negatives:def}{nombre relatif} n'est jamais négatif, et en déduire
        qu'aucun \omterm{def:g7:negatives:def}{nombre relatif} $x$ ne vérifie $x^2 = -4$.
  \item Zoé affirme~: «~$(-1)^{2026} + (-1)^{2027} = 0$, et plus
        généralement deux puissances consécutives de $-1$ s'annulent
        toujours.~» A-t-elle raison~? Justifier.
\end{enumerate}

\medskip
\noindent\textbf{Partie III --- La \omterm{def:g1:addition:def}{somme} alternée.}
Les règles des signes une fois assurées, on peut calculer des \omterm{def:g1:addition:def}{sommes}
mêlant les deux signes. Pour un entier $n \geq 1$, notons
\[
A_n = 1 - 2 + 3 - 4 + \dots
\]
la \omterm{def:g1:addition:def}{somme} des entiers de $1$ à $n$ affectés de signes alternés~: le
dernier terme vaut $+n$ lorsque $n$ est \omterm{def:g3:numbers:evenodd}{impair}, et $-n$ lorsque $n$
est \omterm{def:g3:numbers:evenodd}{pair}. Ainsi $A_1 = 1$ et $A_2 = 1 - 2 = -1$.
\begin{enumerate}[resume]
  \item Calculer $A_3$, $A_4$, $A_5$, $A_6$ et $A_7$. Que
        conjecture-t-on~?
  \item Supposons $n$ \omterm{def:g3:numbers:evenodd}{pair}. Regrouper les termes deux par deux,
        $(1 - 2) + (3 - 4) + \dots$, compter les paires, et démontrer
        que $A_n = -\frac{n}{2}$.
  \item Supposons $n$ \omterm{def:g3:numbers:evenodd}{impair}. Expliquer pourquoi
        $A_n = A_{n-1} + n$, et déduire de la question précédente que
        $A_n = \frac{n + 1}{2}$.
  \item Calculer $A_{100}$, $A_{101}$ et $A_{2026}$.
  \item Alma a calculé $A_{101} - A_{100} = 51 - (-50) = 101$ et a
        été surprise de trouver exactement $101$. Aurait-elle dû
        l'être~? Expliquer en une phrase, et donner la valeur de
        $A_{2027} - A_{2026}$ sans aucun calcul supplémentaire.
\end{enumerate}
\end{problem}
